\input{estilo}
% El número de arriba dice qué es (el último timestamp que le llegó
% a cada réplica); el líder va rotulado abajo y no con una L junto al número, como
% en la pizarra; y
% a la réplica atrasada se le ve el hueco de (x, 11), que viaja por la flecha de Paxos.
\begin{document}
\begin{figura}\begin{tikzpicture}
  \tikzset{replica/.style={draw=tinta, fill=relleno, rounded corners=1.5pt, minimum width=20mm,
    minimum height=17mm, align=center, anchor=north}}
  \node[replica] (a) at (0,0) {(\var{x}, 10)\\[2pt](\var{x}, 11)};
  \node[replica, double=relleno, double distance=1.2pt] (l) at (3.2,0) {(\var{x}, 10)\\[2pt](\var{x}, 11)};
  \node[replica] (c) at (6.4,0) {};
  \node at ($(c.center)+(0,0.22)$) {(\var{x}, 10)};
  \node[draw=suave, densely dashed, text=suave, inner sep=1.5pt] at ($(c.center)+(0,-0.26)$) {(\var{x}, 11)};
  \foreach \n/\t in {a/11, l/11, c/10} \node[anchor=south] at (\n.north) {\t};
  \node[etiq, text=suave, anchor=east, align=right] at ($(a.north west)+(-0.15,0.25)$) {último\\timestamp};
  \node[etiq, anchor=north] at ($(l.south)+(0,-0.08)$) {líder};
  \node[etiq, text=suave, anchor=north] at ($(c.south)+(0,-0.08)$) {todavía no llegó};
  \draw[msg] (l.west) -- (a.east);
  \draw[msg] (l.east) -- (c.west) node[midway, above, etiq] {(\var{x}, 11)};
  \node[anchor=west] (tx) at (8.2,-0.15) {Tx R/O @ \underline{12}};
  \draw[->] (8.3,-1.05) -- ($(c.east)+(0,-0.25)$);
  \node[text=acento, anchor=west, align=left] at (8.4,-1.25) {lee valor viejo\\(\var{x}, 10)};
\end{tikzpicture}\end{figura}
\end{document}
